\documentclass[12pt,a4paper]{article}
\usepackage[utf8]{inputenc}
\usepackage[T2A]{fontenc}
\usepackage[russian]{babel}
\usepackage{amsmath,amssymb}
\usepackage{booktabs,array}
\usepackage{geometry}
\usepackage{microtype}
\geometry{margin=2.3cm}

\title{Shared-holonomy two-sector gate\\
для family-selector и tensor-Hessian}
\author{}
\date{5 августа 2026 г.}

\begin{document}
\maketitle

\section{Проверяемая гипотеза}

Четыре состояния spin-меню несут permutation-представление \(S_4\), а
семейный триплет является его трёхмерным подпространством с нулевой суммой
компонент. Проверяется минимальная объединяющая гипотеза: один и тот же
постоянный holonomy \(g\in S_4\), возможно умноженный на центральный знак,
одновременно
\begin{enumerate}
    \item задаёт outside-\(D_8\) incidence-оператор, расширяющий family-алгебру
    до \(M_3\);
    \item фиксирует три boundary phase tensor-triplet на \(S^1\) и даёт
    полный отклик \(H=\pi^4\).
\end{enumerate}

Для фазы \(a\ne0\) используется точная сумма
\begin{equation}
 \sum_{n\in\mathbb Z}\frac1{(n+a)^4}
 =\pi^4\frac{2+\cos(2\pi a)}{3\sin^4(\pi a)},
\end{equation}
а для периодического канала удаляется нулевая мода:
\begin{equation}
 \sum_{n\in\mathbb Z\setminus\{0\}}\frac1{n^4}=\frac{\pi^4}{45}.
\end{equation}

\section{Полный перебор}

Перебраны все \(24\) элемента \(S_4\) и оба центральных знака. Результаты
зависят только от cycle type:
\begin{center}
\begin{tabular}{cccc}
\toprule
Тип & Знак & Фазы на триплете & \(H/\pi^4\) \\
\midrule
(1+1+1+1) & (+) & (0,0,0) & (1/15) \\
(1+1+1+1) & (-) & (1/2,1/2,1/2) & (1) \\
(2+1+1) & (+) & (0,0,1/2) & (17/45) \\
(2+1+1) & (-) & (0,1/2,1/2) & (31/45) \\
(2+2) & (+) & (0,1/2,1/2) & (31/45) \\
(2+2) & (-) & (0,0,1/2) & (17/45) \\
(3+1) & (+) & (0,1/3,2/3) & (9/5) \\
(3+1) & (-) & (1/4,1/2,3/4) & (17/3) \\
(4) & (+) & (1/4,1/2,3/4) & (17/3) \\
(4) & (-) & (0,1/4,3/4) & (241/45) \\
\bottomrule
\end{tabular}
\end{center}

Ровно один из \(48\) знаковых вариантов даёт \(H=\pi^4\): центральный
antiperiodic holonomy (-I). Но его family-действие скалярно и не расширяет
исходную редуцируемую алгебру. Напротив, все двенадцать outside-\(D_8\)
incidence-направлений типов transposition и three-cycle, дающих полную
\(M_3\)-алгебру, промахиваются по tensor-отклику. Ближайший случай имеет
\(H=31\pi^4/45\), то есть относительный промах \(14/45\).

\section{Вердикт}

\begin{equation}
 \boxed{
 \text{один постоянный signed }S_4\text{-holonomy}
 \quad\not\Rightarrow\quad
 \bigl(M_3\text{ family selector}\bigr)
 \land
 \bigl(H=\pi^4\bigr)
 }
\end{equation}

Это полезный отрицательный результат: впервые оба operator seed проверены в
одной parameter-free конструкции. Простая идентификация их через один group
element закрыта. Следующий допустимый класс должен быть богаче: non-Abelian
boundary connection с path ordering, Wilson-оператор на нескольких рёбрах или
один parent symmetry principle, порождающий разные, но связанные
представления в family- и tensor-секторах.

Машинный перебор сохранён в
\texttt{s2t\_shared\_holonomy\_two\_sector\_audit.py} и
\texttt{s2t\_shared\_holonomy\_two\_sector\_results.json}.

\end{document}